\expandafter\ifx\csname CAPETINCLUDEMODE\endcsname\relax
\documentclass[11pt,a4paper]{article}
% preamble.tex — packages et commandes communes a toutes les fiches CAPET SII IE
\input{glyphtounicode}
\pdfgentounicode=1
\usepackage{cmap}
\usepackage[T1]{fontenc}
\usepackage[utf8]{inputenc}
\usepackage[french]{babel}
\usepackage[a4paper, top=1.6cm, bottom=1.6cm, left=1.4cm, right=1.4cm]{geometry}
\setlength{\headheight}{14.5pt}
\usepackage{amsmath, amssymb, amsfonts}
\usepackage{siunitx}
\sisetup{locale=FR, output-decimal-marker={,}, inter-unit-product=\ensuremath{{}\cdot{}}}
\usepackage{array, booktabs, multirow, tabularx}
\usepackage{xcolor}
\usepackage[most, breakable]{tcolorbox}
\usepackage{multicol}
\usepackage{enumitem}
\usepackage{mathtools}
\usepackage{lmodern}
\usepackage{microtype}
\usepackage{fancyhdr}
\usepackage{tikz}
\usetikzlibrary{positioning,arrows.meta}
\usepackage{lastpage}

% ── Couleurs par domaine ──────────────────────────────────────────────
\definecolor{colElec}{RGB}{26, 82, 148}       % Électrotechnique — bleu foncé
\definecolor{colElecLight}{RGB}{210, 225, 245} % fond bleu clair
\definecolor{colAuto}{RGB}{34, 120, 66}        % Automatique — vert
\definecolor{colMeca}{RGB}{185, 90, 20}        % Mécanique — orange
\definecolor{colTherm}{RGB}{160, 40, 40}       % Thermique — rouge
\definecolor{colQT}{RGB}{232, 244, 233}        % fond question type (vert très léger)
\definecolor{colJury}{RGB}{255, 243, 205}      % fond jury (jaune ambre)
\definecolor{colJuryFrame}{RGB}{200, 130, 0}   % bordure jury
\definecolor{colGrayLight}{RGB}{245, 245, 247} % fond gris neutre
\definecolor{colSeparator}{RGB}{180, 180, 180} % gris séparateur

% ── Style tcolorbox de section principale ────────────────────────────
\tcbset{
  ficheSection/.style={
    breakable,
    enhanced,
    colback=#1!7!white,
    colframe=#1,
    colbacktitle=#1,
    coltitle=white,
    fonttitle=\bfseries\normalsize,
    left=4pt, right=4pt, top=3pt, bottom=4pt,
    before skip=6pt, after skip=6pt,
  }
}

% ── Style encadré question type ──────────────────────────────────────
\tcbset{
  questiontype/.style={
    breakable,
    enhanced,
    colback=colQT,
    colframe=colElec!55!black,
    colbacktitle=colElec!55!black,
    coltitle=white,
    fonttitle=\bfseries\small,
    left=4pt, right=4pt, top=3pt, bottom=3pt,
    before skip=5pt, after skip=5pt,
  }
}

% ── Style encadré jury ───────────────────────────────────────────────
\tcbset{
  jury/.style={
    breakable,
    enhanced,
    colback=colJury,
    colframe=colJuryFrame,
    fonttitle=\bfseries\normalsize,
    left=4pt, right=4pt, top=3pt, bottom=3pt,
  }
}

% ── Mise en page ─────────────────────────────────────────────────────
\pagestyle{fancy}
\fancyhf{}
\renewcommand{\headrulewidth}{0.4pt}
\renewcommand{\footrulewidth}{0.3pt}

\setlist[itemize]{leftmargin=1.2em, itemsep=1pt, topsep=2pt, parsep=0pt}
\setlist[enumerate]{leftmargin=1.2em, itemsep=3pt, topsep=2pt, parsep=0pt}

\setlength{\parskip}{2pt}
\setlength{\columnsep}{10pt}

\def\CAPETBEGINFILE{\begin{document}}
\def\CAPETENDFILE{\end{document}}
\else
\def\CAPETBEGINFILE{}
\def\CAPETENDFILE{}
\fi

\fancyhead[L]{\textbf{\textcolor{colElec}{T20}} - Filtres RC/RL - 1\textsuperscript{er} et 2\textsuperscript{e} ordre}
\fancyhead[R]{\small\textit{CAPET SII IE - Session 2026}}
\fancyfoot[C]{\small\thepage\,/\,\pageref{LastPage}}

\CAPETBEGINFILE

% -----------------------------------------------------------------------
%  BLOC EN-TeTE
% -----------------------------------------------------------------------

\begin{tcolorbox}[
  enhanced, breakable,
  colback=colElec,
  colframe=colElec,
  left=6pt, right=6pt, top=5pt, bottom=5pt,
  before skip=0pt, after skip=4pt
]
  {\color{white}\bfseries\large T20 \quad Filtres RC/RL - 1\textsuperscript{er} et 2\textsuperscript{e} ordre}\\[2pt]
  {\color{colElec!30!white}\small
    \textbf{Domaine :} electronique - Traitement analogique du signal \quad
    \textbf{Presence :} \textbf{7/13} sujets (2013-2026) - theme \emph{tres frequent}
  }\\[1pt]
  {\color{colElec!20!white}\footnotesize
    Sous-themes : filtre passe-bas RC - filtre RL passe-bas - diagramme de Bode - gabarit frequentiel - filtre du 2\textsuperscript{e} ordre (cascade, RLC) - filtre anti-repliement - Shannon
  }
\end{tcolorbox}

% -----------------------------------------------------------------------
%  1. RAPPELS DE COURS
% -----------------------------------------------------------------------

\begin{tcolorbox}[
  ficheSection=colElec,
  title={1. Rappels de cours},
]

% -- 1.1 Filtre passe-bas RC du 1er ordre ------------------------------
\textbf{1.1 Filtre passe-bas RC - forme canonique}

\smallskip
Schema : $R$ en serie, $C$ en parallele avec la sortie.
\[
  H(j\omega) = \frac{1}{1 + j\,\dfrac{\omega}{\omega_0}}
  \qquad \text{avec} \qquad \omega_0 = \frac{1}{RC}
\]
\begin{center}
\renewcommand{\arraystretch}{1.3}
\begin{tabular}{lll}
\toprule
\textbf{Grandeur} & \textbf{Expression} & \textbf{ $\omega = \omega_0$} \\
\midrule
Module $|H|$ & $\dfrac{1}{\sqrt{1+(\omega/\omega_0)^2}}$ & $\dfrac{1}{\sqrt{2}} \approx 0{,}707$ \\[6pt]
Gain en dB & $-20\log\sqrt{1+(\omega/\omega_0)^2}$ & $-3\,\mathrm{dB}$ \\[4pt]
Phase $\varphi$ & $-\arctan(\omega/\omega_0)$ & $-45^\circ$ \\[4pt]
Temps de reponse - 5\,\% & \multicolumn{2}{l}{$t_{r5\,\%} \approx 3\tau = 3\,RC$} \\
\bottomrule
\end{tabular}
\end{center}

\medskip\hrule\vspace{4pt}

% -- 1.2 Diagramme de Bode asymptotique --------------------------------
\textbf{1.2 Diagramme de Bode - regles asymptotiques}

\smallskip
\begin{center}
\renewcommand{\arraystretch}{1.3}
\begin{tabular}{lll}
\toprule
\textbf{Zone} & \textbf{Gain asymptotique} & \textbf{Phase asymptotique} \\
\midrule
$\omega \ll \omega_0$ & $0\,\mathrm{dB}$ (plateau) & $0^\circ$ \\
$\omega = \omega_0$ & $-3\,\mathrm{dB}$ (reel) & $-45^\circ$ (reel) \\
$\omega \gg \omega_0$ & pente $-20\,\mathrm{dB/decade}$ & $-90^\circ$ \\
\bottomrule
\end{tabular}
\end{center}

\medskip\hrule\vspace{4pt}

% -- 1.3 Filtre RL passe-bas -------------------------------------------
\textbf{1.3 Filtre RL passe-bas}

\smallskip
$R$ en parallele avec $L$ : meme forme que RC avec $\tau = L/R$.
\[
  H(j\omega) = \frac{1}{1 + j\,\dfrac{\omega}{\omega_0}}
  \qquad \omega_0 = \frac{R}{L}
\]
Les enroulements d'une machine asynchrone se comportent comme un filtre RL passe-bas naturel qui attenue les harmoniques de la MLI.

\medskip\hrule\vspace{4pt}

% -- 1.4 Filtre du 2e ordre --------------------------------------------
\textbf{1.4 Filtre passe-bas du 2\textsuperscript{e} ordre - forme canonique}

\[
  H(j\omega) = \frac{H_0}{1 + 2m\,\dfrac{j\omega}{\omega_0} + \left(\dfrac{j\omega}{\omega_0}\right)^2}
\]

\begin{center}
\renewcommand{\arraystretch}{1.3}
\begin{tabular}{lll}
\toprule
\textbf{Parametre} & \textbf{Signification} & \textbf{Effet} \\
\midrule
$\omega_0$ & Pulsation propre (rad/s) & Position de la coupure \\
$m$ & Facteur d'amortissement & Forme de la reponse \\
$Q = 1/(2m)$ & Facteur de qualite & Selectivite \\
\midrule
$m > 1$ & Sur-amorti & Pas d'oscillations \\
$m = 1$ & Critique & Retour le plus rapide sans depassement \\
$m < 1$ & Sous-amorti & Resonance, pic - $\omega_r = \omega_0\sqrt{1-2m^2}$ \\
\bottomrule
\end{tabular}
\end{center}

\medskip\hrule\vspace{4pt}

% -- 1.5 Cascade de 2 filtres 1er ordre --------------------------------
\textbf{1.5 Cascade de deux filtres RC identiques}

\[
  H(j\omega) = \left(\frac{1}{1 + j\omega/\omega_0}\right)^2
\]
Pente asymptotique $-40\,\mathrm{dB/decade}$ ; amortissement $m = 1$ (pas d'oscillations ni de resonance).

\medskip\hrule\vspace{4pt}

% -- 1.6 Filtre anti-repliement ----------------------------------------
\textbf{1.6 Filtre anti-repliement (Shannon)}

Place avant le CAN pour eliminer les composantes - $f > f_e/2$ :
\[
  f_e \geq 2\,f_{max} \qquad \Rightarrow \qquad f_c \leq \frac{f_e}{2}
\]

\end{tcolorbox}

% -----------------------------------------------------------------------
%  2. FORMULAIRE SYNTHeTIQUE
% -----------------------------------------------------------------------

\begin{tcolorbox}[
  ficheSection=colElec,
  title={2. Formulaire synthetique},
]
\small
\begin{multicols}{2}
\renewcommand{\arraystretch}{1.45}
\begin{tabular}{@{}p{4.8cm}ll@{}}
\toprule
\textbf{Formule} & \textbf{Usage} & \textbf{U.} \\
\midrule
\multicolumn{3}{@{}l}{\footnotesize\textit{Filtre 1\textsuperscript{er} ordre RC}} \\
$\omega_0 = 1/(RC)$ & Coupure & rad/s \\
$|H(\omega_0)| = 1/\sqrt{2}$ & Module - $\omega_0$ & - \\
$G_{dB}(\omega_0) = -3\,\mathrm{dB}$ & Gain & dB \\
$\varphi(\omega_0) = -45^\circ$ & Phase & - \\
$t_{r5\%} \approx 3\,RC$ & Transitoire & s \\
\midrule
\multicolumn{3}{@{}l}{\footnotesize\textit{Filtre RL}} \\
$\omega_0 = R/L$ & Coupure & rad/s \\
$\tau = L/R$ & Cste de temps & s \\
\bottomrule
\end{tabular}

\columnbreak

\renewcommand{\arraystretch}{1.45}
\begin{tabular}{@{}p{4.8cm}ll@{}}
\toprule
\textbf{Formule} & \textbf{Usage} & \textbf{U.} \\
\midrule
\multicolumn{3}{@{}l}{\footnotesize\textit{Filtre 2\textsuperscript{e} ordre}} \\
$H = H_0 / (1 + 2m\,(j\omega/\omega_0) + (j\omega/\omega_0)^2)$ & Forme canon. & - \\
$Q = 1/(2m)$ & Qualite & - \\
$\omega_r = \omega_0\sqrt{1-2m^2}$ & Resonance ($m<1/\sqrt{2}$) & rad/s \\
Pente HF : $-40\,\mathrm{dB/dec}$ & Bode & - \\
\midrule
\multicolumn{3}{@{}l}{\footnotesize\textit{Shannon / anti-repliement}} \\
$f_e \geq 2\,f_{max}$ & Shannon & Hz \\
$f_c \leq f_e/2$ & Filtre AR & Hz \\
\midrule
\multicolumn{3}{@{}l}{\footnotesize\textit{Gabarit}} \\
Attenuation $>100$ ($>40\,\mathrm{dB}$) $\Rightarrow$ 2\textsuperscript{e} ordre & & \\
\bottomrule
\end{tabular}
\end{multicols}
\normalsize
\end{tcolorbox}

% -----------------------------------------------------------------------
%  3. QUESTIONS TYPES
% -----------------------------------------------------------------------

\begin{tcolorbox}[
  ficheSection=colElec,
  title={3. Questions types - extraites des sujets reels},
]

% -- QT1 -------------------------------------------------------------

% Mini-procedure de resolution type copie concours
\begin{tcolorbox}[questiontype,title={Procedure type copie concours (5 etapes)}]
\begin{enumerate}[leftmargin=1.5em]
  \item \textbf{Qualifier le modele} : hypotheses, domaine de validite, grandeurs d'entree/sortie.
  \item \textbf{Ecrire les lois} : equations de base avec symboles et unites SI.
  \item \textbf{Isoler l'inconnue} : expression litterale, puis application numerique tracee.
  \item \textbf{Controler l'ordre de grandeur} : coherence physique et unite finale.
  \item \textbf{Conclure concours} : phrase courte orientee decision technique.
\end{enumerate}
\end{tcolorbox}

% Visuel de methode : chemin de resolution CAPET
\begin{center}
\begin{tikzpicture}[font=\small, node distance=7mm and 6mm, >=Stealth]
  \node[draw, rounded corners, fill=colElec!8, inner sep=4pt] (n1) {Notions};
  \node[draw, rounded corners, fill=colElec!8, right=of n1, inner sep=4pt] (n2) {Methode};
  \node[draw, rounded corners, fill=colElec!8, right=of n2, inner sep=4pt] (n3) {Application annale};
  \node[draw, rounded corners, fill=colElec!8, right=of n3, inner sep=4pt] (n4) {Verification jury};
  \draw[->, thick, colElec] (n1) -- (n2);
  \draw[->, thick, colElec] (n2) -- (n3);
  \draw[->, thick, colElec] (n3) -- (n4);
\end{tikzpicture}
\end{center}
\begin{tcolorbox}[questiontype,
  title={QT1 - Filtre RC - fonction de transfert et diagramme de Bode \hfill\normalfont\footnotesize E1 2026 (B3-B4)}
]
\textbf{enonce.}
Filtre passe-bas RC avec $R = \SI{1}{\kilo\ohm}$ et $C = \SI{100}{\nano\farad}$. Montrer que $H_{F1}(j\omega) = 1/(1+j\omega/\omega_0)$. Exprimer $\omega_0$. Tracer le diagramme de Bode asymptotique.

\tcbline

\textbf{Corrige.}
Par diviseur de tension ($Z_C = 1/(jC\omega)$) :
\[
  H = \frac{Z_C}{R + Z_C} = \frac{1/(jC\omega)}{R + 1/(jC\omega)} = \frac{1}{1 + jRC\omega}
  \qquad \Rightarrow \qquad \omega_0 = \frac{1}{RC} = \frac{1}{10^3 \times 10^{-7}} = \SI{10000}{\radian\per\second}
\]
Bode : plateau $0\,\mathrm{dB}$ pour $\omega \ll \omega_0$ ; pente $-20\,\mathrm{dB/decade}$ pour $\omega \gg \omega_0$ ; point reel - $\omega_0$ : $-3\,\mathrm{dB}$, $-45^\circ$.
\end{tcolorbox}

\vspace{4pt}

% -- QT2 -------------------------------------------------------------

\begin{tcolorbox}[questiontype,
  title={QT2 - Lecture graphique sur le Bode - calcul de $v_3(t)$ \hfill\normalfont\footnotesize E1 2026 (B5-B6)}
]
\textbf{enonce.}
Signal d'entree $v_2(t) = 2 + 0{,}5\sin(2\pi \cdot 900 \cdot t)$. Avec $f_0 = \SI{1592}{\hertz}$ ($\omega_0 = 10^4\,\mathrm{rad/s}$), determiner le signal de sortie $v_3(t)$ et estimer le temps de reponse - 5\,\%.

\tcbline

\textbf{Corrige.}

Composante DC ($\omega = 0$) : $|H(0)| = 1$ - $V_{DC} = 2\,\mathrm{V}$.

Composante AC ($f_2 = 900\,\mathrm{Hz}$ - $\omega_2 = 2\pi \times 900 \approx 5655\,\mathrm{rad/s}$) :
\[
  |H(\omega_2)| = \frac{1}{\sqrt{1+(5655/10000)^2}} = \frac{1}{\sqrt{1+0{,}32}} \approx 0{,}870
  \qquad
  V_{AC,out} = 0{,}5 \times 0{,}870 \approx 0{,}435\,\mathrm{V}
\]
\[
  v_3(t) \approx 2 + 0{,}435\sin(2\pi \times 900\,t + \varphi) \,\mathrm{V}
\]

Temps de reponse : $t_{r5\%} \approx 3\,RC = 3 \times 10^3 \times 10^{-7} = 0{,}3\,\mathrm{ms}$ - cahier des charges ($< 50\,\mathrm{ms}$) satisfait, mais l'attenuation - 900 Hz est insuffisante (filtre du 1\textsuperscript{er} ordre) - filtre du 2\textsuperscript{e} ordre necessaire.
\end{tcolorbox}

\vspace{4pt}

% -- QT3 -------------------------------------------------------------

\begin{tcolorbox}[questiontype,
  title={QT3 - Filtre du 2\textsuperscript{e} ordre - choix des composants et reponse indicielle \hfill\normalfont\footnotesize E1 2026 (B9)}
]
\textbf{enonce.}
On souhaite attenuer $> 40\,\mathrm{dB}$ - $\omega_2 = 5655\,\mathrm{rad/s}$ avec un temps de reponse - 5\,\% $< 50\,\mathrm{ms}$. Proposer $R$ et $C$ pour un filtre 2\textsuperscript{e} ordre (cascade de deux RC). Justifier si la reponse indicielle presente des oscillations.

\tcbline

\textbf{Corrige.}

Condition de gain : $|H(\omega_2)|^2 = [1/(1+(\omega_2/\omega_0)^2)] \leq 1/100$ - $\omega_0 \leq \omega_2/\sqrt{99} \approx 568\,\mathrm{rad/s}$ - $f_0 \leq 90\,\mathrm{Hz}$.

Condition de rapidite : $t_{r5\%} \approx 3RC \leq 50\,\mathrm{ms}$ - $RC \leq 16{,}7\,\mathrm{ms}$.

Choix : $RC = 1/(568) \approx 1{,}76\,\mathrm{ms}$ - par ex. $R = \SI{17{,}6}{\kilo\ohm}$, $C = \SI{100}{\nano\farad}$.

La cascade de deux filtres RC identiques donne $m = 1$ (regime critique) - \textbf{pas d'oscillations}.
\end{tcolorbox}

\end{tcolorbox}

% -----------------------------------------------------------------------
%  4. ATTENTES DU JURY
% -----------------------------------------------------------------------

\begin{tcolorbox}[
  jury,
  title={4. Attentes du jury - Erreurs recurrentes et points de vigilance},
]

\begin{enumerate}[label=\textbf{\arabic*.}, leftmargin=*, itemsep=5pt]

  \item \textbf{Maetriser la forme canonique} $H = 1/(1+j\omega/\omega_0)$ et l'identifier dans tout schema RC ou RL - ne pas recalculer depuis les equations de circuit si la forme est deje fournie.

  \item \textbf{Ne pas confondre $\omega_0$ et $f_0$} : $\omega_0 = 2\pi f_0$ en rad/s, $f_0$ en Hz. Donner les unites explicitement dans chaque resultat.

  \item \textbf{Gain - $\omega_0$ vaut $-3\,\mathrm{dB}$, pas 0\,dB.} Erreur frequente : les candidats ecrivent $|H(\omega_0)| = 0$ au lieu de $1/\sqrt{2}$.

  \item \textbf{Justifier l'ordre du filtre depuis le gabarit} : pente disponible par ordre = $-20\,\mathrm{dB/decade}$. Pour attenuer $>40\,\mathrm{dB}$ - 2\textsuperscript{e} ordre minimum. Ne pas choisir l'ordre par intuition.

  \item \textbf{Filtre anti-repliement} : systematiquement mentionne des qu'un CAN apparaet dans le schema. Rappeler le theoreme de Shannon, puis deduire $f_c \leq f_e/2$.

  \item \textbf{Cascade - RLC} : deux filtres RC en cascade donnent $m = 1$ (pas de resonance), contrairement - un circuit RLC serie ou parallele.

\end{enumerate}

\medskip
\begin{tcolorbox}[
  colback=colJuryFrame!12!colJury,
  colframe=colJuryFrame,
  left=4pt, right=4pt, top=2pt, bottom=2pt,
  fontupper=\small\itshape
]
Rapport jury 2026 : "La partie B (filtre, Bode) a ete abordee par la majorite des candidats mais les traces asymptotiques sont souvent incorrects (pentes, phases, positionnement de $\omega_0$). Le jury attend le trace reel esquisse \emph{en plus} des asymptotes."\\
Rapport jury 2024 : "Les enroulements de la machine asynchrone constituent un filtre passe-bas RL naturel - les candidats doivent connaetre ce resultat et savoir le justifier."
\end{tcolorbox}

\end{tcolorbox}

% -----------------------------------------------------------------------
%  PIED - PReSENCE DANS LES SUJETS
% -----------------------------------------------------------------------

\begin{tcolorbox}[
  enhanced, breakable,
  colback=colGrayLight,
  colframe=colSeparator,
  left=4pt, right=4pt, top=3pt, bottom=3pt,
  before skip=4pt, after skip=0pt,
  fonttitle=\small\bfseries,
  title={Presence dans les sujets analyses (2013-2026)},
]
\begin{center}
\footnotesize\renewcommand{\arraystretch}{1.15}
\begin{tabular}{lccccccccccccc}
\toprule
\textbf{Annee} & 13 & 14 & 15 & 16 & 17 & 18 & 19 & 20 & 21 & 22 & 23 & 24 & 25 \\
\midrule
E1 (commune)   & - & - & - & - & - & - & - & - & - & - & $\checkmark$ & - & $\checkmark$ \\
E2 discipl. IE & - & - & - & - & - & $\checkmark$ & - & $\checkmark$ & - & - & - & $\checkmark$ & - \\
E2 appl. comm. & - & - & - & - & - & - & - & - & - & - & - & - & - \\
\bottomrule
\end{tabular}
\end{center}
\smallskip
\footnotesize\textbf{Sous-themes frequents :}
filtre passe-bas RC ($\omega_0$, Bode) - diagramme de Bode asymptotique + reel - gabarit frequentiel - filtre anti-repliement + Shannon - cascade de 2 filtres RC

\medskip
\textbf{Ordres de grandeur (reperes rapides).}
\begin{itemize}[leftmargin=1.3em]
  \item Reseau BT FR : tension phase-neutre typique \SI{230}{\volt}, tension phase-phase typique \SI{400}{\volt}, frequence \SI{50}{\hertz}.
  \item Courants d'etude frequents en sujet : de quelques \si{\ampere} a quelques dizaines de \si{\ampere}.
  \item En conversion d'energie, un rendement attendu reste generalement entre 0{,}8 et 0{,}98 selon la technologie.
\end{itemize}

\medskip
\textbf{Synthese express (a refaire sans support).}
\begin{enumerate}[leftmargin=1.4em]
  \item Poser les hypotheses du modele et les conventions de signe.
  \item Ecrire 2 a 4 relations cles avec unites.
  \item Realiser un calcul numerique complet et verifier l'ordre de grandeur.
  \item Justifier le resultat par un argument technique de niveau CAPET.
\end{enumerate}

\end{tcolorbox}

\CAPETENDFILE









